% !TEX root = ../main.tex
\section{模型假设}

结合题目条件和前两问的建模需要，作如下假设：

\begin{enumerate}
  \item 测向机返回的示向度均已按照题目规定归一化至 $[0^\circ,360^\circ)$；角度误差按圆周最短角距离计算，即
  $\left|\operatorname{wrap}_{(-\pi,\pi]}(\theta_{\mathrm{obs}}-\theta_{\mathrm{true}})\right|\leq1^\circ$；
  \item 前后两次观测针对同一频道、同一干扰源，且干扰源在定位过程中保持静止；
  \item 忽略测向机与干扰源之间的高度差异，将检测点和干扰源视为同一平面上的点；
  \item 不为测量误差指定额外的概率分布，而将其作为区间 $[-1^\circ,1^\circ]$ 内的有界不确定量处理；
  \item 第二问中的干扰源为全向干扰源，在其有效接收半径内不存在由发射方向造成的信号盲区；
  \item 数值容差仅用于计算，不作为额外物理误差；方向、距离与叉积判定采用相应量纲容差。
\end{enumerate}
